% !TeX root = ../../Book4.tex
% 纲要标题：### 6.1 标记与等价关系
% 纲要位置：几何学计划纲要.md，第 3683 行。
% 正文状态：未撰写；本文件目前是框架。
% 本次扩充的范围、先修与边界：先修为 I.6 的闭曲面分类、II.2 的一阶同调及闭曲面的同调交叉配对，并使用本节已有的标记、同位和群作用。固定闭连通定向曲面，构造 Dehn 扭转并核对方向和交叉数符号，计算其同调作用，将 Torelli 群定义为作用的核。Nielsen–Thurston 三分法限亏格至少二的精确陈述与代表模型；不承担完整分类证明、映射类群的全部生成关系或 Torelli 群的专门结构理论。

% 纲要内容开始
% 
% > 补充写作范围与外部输入：先修为 I.6 的闭曲面分类、II.2 的一阶同调及闭曲面的同调交叉配对，并使用本节已有的标记、同位和群作用。固定闭连通定向曲面，构造 Dehn 扭转并核对方向和交叉数符号，计算其同调作用，将 Torelli 群定义为作用的核。Nielsen–Thurston 三分法限亏格至少二的精确陈述与代表模型；不承担完整分类证明、映射类群的全部生成关系或 Torelli 群的专门结构理论。
% 
% * 标记曲面、同位和双全纯
% * Teichmüller 空间与模空间的区别
% * 映射类群作用及稳定子
% 
% * Dehn 扭转的环带构造与同位类
% * Dehn 扭转对一阶同调的作用
% * Torelli 群作为同调作用的核
% * Nielsen–Thurston 三分法的陈述与模型
% 
% 
% 纲要内容结束

\section{标记与等价关系}
\label{g4:sec:0601}

待撰写

\subsection{标记曲面、同位和双全纯}
\label{g4:subsec:0601:01}
待撰写

\subsection{Teichmüller 空间与模空间的区别}
\label{g4:subsec:0601:02}
待撰写

\subsection{映射类群作用及稳定子}
\label{g4:subsec:0601:03}
待撰写

\subsection{Dehn 扭转的环带构造与同位类}
\label{g4:subsec:0601:04}
待撰写

\subsection{Dehn 扭转对一阶同调的作用}
\label{g4:subsec:0601:05}
待撰写

\subsection{Torelli 群作为同调作用的核}
\label{g4:subsec:0601:06}
待撰写

\subsection{Nielsen–Thurston 三分法的陈述与模型}
\label{g4:subsec:0601:07}
待撰写

% 纲要条目排为编号小节标题，并进入目录与书签；说明段仍保留于纲要注释中。
